\documentclass[a4paper]{article}
\usepackage{amsfonts}
\usepackage{amsmath}
\usepackage{amssymb}
\usepackage{graphicx}     
\begin{document}
  
\noindent \large Auteur: Viki Lauvrys        \\
\noindent \large Studierichting: Informatica \\
\noindent \large Oefeningenreeks 1C nr. 46.2 \\

\medskip

\normalsize

$ z^4 = -1 = 1 \cdot \operatorname{cis}(\pi) $\\

$ \Rightarrow z_k = \operatorname{cis}\left( \dfrac{\pi + 2k\pi}{4} \right) $ met $ k \in \mathbb{Z} $ \\

De vierdemachtswortels zijn dan: \\

$ \Rightarrow z_0 = \operatorname{cis}\left( \dfrac{\pi}{4} \right) = \dfrac{\sqrt{2}}{2} + \dfrac{\sqrt{2}i}{2} $  \\

$ \Rightarrow z_1 = \operatorname{cis}\left( \dfrac{3\pi}{4} \right) = -\dfrac{\sqrt{2}}{2} + \dfrac{\sqrt{2}i}{2} $  \\

$ \Rightarrow z_2 = \operatorname{cis}\left( \dfrac{5\pi}{4} \right) = -\dfrac{\sqrt{2}}{2} - \dfrac{\sqrt{2}i}{2} $  \\

$ \Rightarrow z_3 = \operatorname{cis}\left( \dfrac{7\pi}{4} \right) = \dfrac{\sqrt{2}}{2} - \dfrac{\sqrt{2}i}{2} $  \\

\begin{thebibliography}{999}
%\bibitem{a} Referentie
\end{thebibliography}
\end{document}